\chapter{Can I Completely Change?}

This chapter follows J. Krishnamurti's inquiry into the possibility of fundamental change. There is no blackboard mathematics in this discussion, and no validated lecture frame contributes equations or diagrams. The mathematical notation below is therefore deliberately modest: it is a compact way of keeping the logic of the dialogue visible. We use it as a map of distinctions, not as a replacement for the spoken inquiry.

\section{The Old Question}

The conversation begins by returning to a question already under way: why do human beings live as they do? The fact is not obscure. There is turmoil, confusion, sorrow, conflict, violence, and a long public history of proposed solutions. Gurus, priests, political systems, books, disciplines, and therapies all offer a way out. The offer has the same form again and again: do this and things will come right.

Yet the fact remains. Human beings know a great deal. Scientific knowledge, biology, sociology, political history, religious history, and psychological explanation are all available. But knowledge has not made human beings live in order, happily, intelligently, without the same chaotic movement repeating itself.

We can write the opening problem as a deliberately simple implication that has failed:

\begin{equation}
\mathrm{knowledge}+\mathrm{method}+\mathrm{authority}
\not\Rightarrow
\mathrm{order}.
\end{equation}

That is not a theorem. It is the fact from which the inquiry begins. If explanations and methods have been tried for centuries and the condition remains, then the question cannot merely be, ``Which method should we try next?'' The question becomes more radical:

\begin{equation}
Q_0:\qquad
\mathrm{What\ is\ the\ root\ of\ human\ disorder?}
\end{equation}

The word ``root'' matters. The lecture does not begin by asking for a better adjustment, a more efficient discipline, or a more consoling belief. It asks why the whole structure continues.

\section{Secondary Explanations}

Several explanations are tested early in the discussion. Perhaps sorrow gives security. Perhaps people get used to conflict and miss it when it is absent. Perhaps war appears to release one from boredom. Perhaps the attraction of war lies in the removal of responsibility: the generals decide, the individual obeys. Perhaps human beings are cynical because history seems to show that nothing can be done.

These explanations are not dismissed because they are foolish. They are dismissed because they are partial. A secondary explanation may describe one local effect without touching the root.

\begin{table}
\centering
\begin{tabular}{p{0.34\linewidth}p{0.54\linewidth}}
\hline
\textbf{Proposed explanation} & \textbf{Why it remains secondary} \\
\hline
Habit & We may become used to conflict, but habit does not explain why the conflict began. \\
False security & Sorrow may feel familiar, but familiarity is not the root of the whole disorder. \\
War as escape & War may interrupt boredom or private emptiness, but the attraction of escape still has to be explained. \\
Transferred responsibility & Obedience may relieve the burden of choice, but it depends on a deeper confusion. \\
Cynicism & The belief that nothing can change may sustain disorder, but it is already a symptom of disorder. \\
\hline
\end{tabular}
\caption{Secondary explanations are useful only if they return us to the root question.}
\end{table}

The inquiry therefore keeps returning to the same pressure:

\begin{equation}
\mathrm{secondary\ cause}
\neq
\mathrm{root}.
\end{equation}

This is the first discipline of the chapter. We do not collect explanations. We ask whether an explanation reaches the centre of the problem.

\section{Fragmentation And The Demand To Change}

The discussion then becomes more direct. Society may be neurotic; but the one who judges society is not outside society. The disorder is not merely ``out there.'' The inquiry turns toward the individual: why do I not change?

At this point the word ``fragmentation'' enters the argument. Human life is divided into political action, religious action, social action, private ambition, pleasure, frustration, belief, and fear. These fragments do not form an intelligent whole. They often oppose one another, and the person is not separate from that division. We are that fragmentation.

\begin{figure}
\centering
\begin{tikzpicture}[
  node distance=1.05cm,
  box/.style={draw, rectangle, align=center, minimum width=0.74\linewidth, minimum height=0.75cm},
  smallbox/.style={draw, rectangle, align=center, minimum width=0.28\linewidth, minimum height=0.75cm}
]
\node[box] (human) {Human life appears as fragmented action};
\node[smallbox, below left=1.0cm and -0.1cm of human] (pol) {political\\action};
\node[smallbox, below=1.0cm of human] (rel) {religious\\action};
\node[smallbox, below right=1.0cm and -0.1cm of human] (soc) {social\\action};
\node[box, below=2.3cm of human] (conflict) {separate movements, often in conflict};
\node[box, below=0.8cm of conflict] (we) {``we'' are not outside this fragmentation};
\draw[->] (human) -- (pol);
\draw[->] (human) -- (rel);
\draw[->] (human) -- (soc);
\draw[->] (pol) -- (conflict);
\draw[->] (rel) -- (conflict);
\draw[->] (soc) -- (conflict);
\draw[->] (conflict) -- (we);
\end{tikzpicture}
\caption{A transcript-backed reconstruction of fragmentation: political, religious, and social action separated from the whole.}
\end{figure}

We can name the condition by a symbol, provided we remember that the symbol is only a convenience. Let \(D\) denote disorder or neurotic confusion, and let \(F\) denote fragmentation. The lecture suggests that the two are not separate:

\begin{equation}
F \subset D,
\qquad
\mathrm{and\ in\ practice}\qquad
\mathrm{me}\in F.
\end{equation}

The person who says, ``I want to change,'' is already speaking from within the field that has to be understood.

\subsection{Question \& Answer}

\paragraph{Question.}
If I want to change, why does change so often reproduce the old?

\paragraph{Answer.}
Because the entity that wants change sets the pattern of change. The pattern may be new, but the planner is old. A new discipline, a new belief, a new book, or a new system may give the feeling of novelty, but the centre using it remains the same centre.

The logic can be written as follows. Let \(I_{\mathrm{old}}\) denote the old centre, and \(P_{\mathrm{new}}\) denote a new pattern or method.

\begin{align}
I_{\mathrm{old}} &\longrightarrow P_{\mathrm{new}},\\
P_{\mathrm{new}} &\longrightarrow \mathrm{changed\ appearance},\\
I_{\mathrm{old}}+P_{\mathrm{new}} &\not\Rightarrow \mathrm{fundamental\ change}.
\end{align}

The last line is the essential one. The pattern changes; the planner remains. Or, in the language of the discussion: the old conquers the new.

\begin{proposition}
A change designed by the old centre cannot by itself be fundamental change.
\end{proposition}

\begin{proof}
The proposed change begins with the old centre: the centre forms an image of what it wants, selects a method, and projects an end. Therefore the structure that is to be changed is already operating as the author of change. The output may have a different colour, name, discipline, or vocabulary, but the source of the movement has not ended. Hence the change remains within the old pattern of becoming.
\end{proof}

\section{Authority Created By Disorder}

The next movement of the lecture rolls the question back into a more concrete form: authority. Who is my authority? The religious teacher, the political leader, the analyst, the system, the book, the discipline, the promise that if I follow now I shall be free later?

The inquiry does not merely reject this or that authority. It asks why authority has such force. The answer is that disorder creates the demand for authority. When I am confused, I want someone else to tell me what to do. When the disorder is intense, the promise of order becomes persuasive.

Let \(A\) denote authority in this psychological sense: not simply public office or technical expertise, but the inward demand that another should direct the movement of my life. Then the lecture's logic is:

\begin{equation}
D \longrightarrow A.
\end{equation}

But if authority is born from disorder, then authority cannot be the root-ending of disorder. The loop is:

\begin{equation}
D \longrightarrow A \longrightarrow D.
\end{equation}

This is one of the cleanest structures in the lecture. Disorder produces the hope that someone else knows. That hope creates obedience, dependence, and imitation. But dependence sustains the disorder that created the authority.

\begin{figure}
\centering
\begin{tikzpicture}[
  node distance=0.85cm,
  box/.style={draw, rectangle, align=center, minimum width=0.76\linewidth, minimum height=0.8cm}
]
\node[box] (d1) {disorder, confusion, fear of not knowing};
\node[box, below=of d1] (need) {need for someone else to tell me what to do};
\node[box, below=of need] (auth) {authority: guru, priest, system, discipline, analyst};
\node[box, below=of auth] (promise) {promise of order through following};
\node[box, below=of promise] (d2) {continued dependence and disorder};
\draw[->] (d1) -- (need);
\draw[->] (need) -- (auth);
\draw[->] (auth) -- (promise);
\draw[->] (promise) -- (d2);
\draw[->] (d2.east) .. controls +(1.2,0) and +(1.2,0) .. (d1.east);
\end{tikzpicture}
\caption{Authority as a loop generated by disorder. This is an editorial reconstruction from the transcript, not a blackboard diagram.}
\end{figure}

\subsection{Question \& Answer}

\paragraph{Question.}
Must I follow authority in order to see my dependence on authority?

\paragraph{Answer.}
The lecture refuses that logic. To say that I must submit to authority in order to discover my dependence on authority is already to strengthen the pattern. Similarly, to say that I must deceive myself in order to understand self-deception is incoherent. If self-deception is genuine, then the deceived mind does not know what it is doing.

The rejected implication is:

\begin{equation}
\mathrm{self\mbox{-}deception}
\not\Rightarrow
\mathrm{understanding\ of\ self\mbox{-}deception}.
\end{equation}

And the parallel implication is:

\begin{equation}
\mathrm{following\ authority}
\not\Rightarrow
\mathrm{freedom\ from\ authority}.
\end{equation}

This is not an anti-intellectual rejection of help. It is more precise. If the root of the appeal to authority is disorder, then the use of authority as the cure leaves the root untouched.

\section{Correct Action Without A Guide}

Once authority is seen in this way, rejection changes its meaning. It is not a reaction, a pose, or a cantankerous refusal. The rejection is sane because the structure has been seen. Authority has been created by disorder; therefore to follow it psychologically is to continue the disorder.

The question then becomes practical. If I reject authority, what is correct action? Who will tell me how to live properly? The communists, Marx, Lenin, Mao, the Pope, the priest, the analyst, the guru? The lecture answers by refusing the premise. If I am neurotic, I cannot find correct action as a formula. And no authority can hand it to me.

So the question shifts:

\begin{align}
Q_1 &: \quad \mathrm{What\ is\ correct\ action?}\\
Q_2 &: \quad \mathrm{Can\ the\ mind\ be\ free\ of\ neuroticism?}
\end{align}

This shift is not evasive. It is exact. Correct action cannot be separated from the state of the mind that acts. If the actor is confused, the action will carry that confusion, however noble its stated aim.

A cautious compression of the argument is:

\begin{equation}
\mathrm{correct\ action}
\quad\hbox{requires}\quad
\mathrm{absence\ of\ inward\ disorder}.
\end{equation}

The lecture does not present this as a moral rule. It presents it as a necessity of observation. If disorder is operating, action springs from disorder.

Here the rejection of authority produces energy. The mind is no longer distributing responsibility among teachers, systems, holy places, or experts. It is no longer looking outward for the next pattern. The seriousness gathers because the question is now one's own life.

\paragraph{A worked logical reduction.}
We can set out the movement in five steps:

\begin{enumerate}
\item Suppose I ask, ``What is the right action?''
\item If I am inwardly confused, my answer will be shaped by confusion.
\item If I ask an authority, the asking itself may come from the same confusion.
\item Therefore the first question cannot be a request for instruction.
\item The first question becomes: can the mind see and be free of its disorder?
\end{enumerate}

This is the lecture's turning point from negation to investigation.

\section{Consciousness, The Word, And Reality}

The inquiry now asks what this disorder is made of. Krishnamurti limits the word consciousness for the moment: consciousness is the collected movement of fragments, memory, conclusions, beliefs, images, words, and thought. The qualification is important. This is not a complete metaphysical definition of consciousness. It is the working field of the present inquiry.

Let \(C\) denote consciousness in this restricted sense. Then:

\begin{equation}
C \approx
\mathrm{beliefs}+\mathrm{memories}+\mathrm{conclusions}
+\mathrm{images}+\mathrm{fragmented\ thought}.
\end{equation}

The ``me'' is not outside this collection. The me is made by the movement of these fragments.

\begin{equation}
\mathrm{me}
\longrightarrow
\mathrm{separation}
\longrightarrow
D.
\end{equation}

The dialogue then encounters a difficulty of language. If I am that movement, how can I look at it? The words themselves may be part of the very system under investigation. We must use words, but we must not take them as fixed things. This brings in one of the central distinctions:

\begin{equation}
\mathrm{word}\neq \mathrm{thing}.
\end{equation}

The word is thought. A belief is thought. A conclusion is thought. But when belief is strong, it appears as reality. To the believer, God, nation, ideology, self, or theory is not merely an idea. It feels real.

The movement can be written:

\begin{align}
\mathrm{word} &\in \mathrm{thought},\\
\mathrm{thought} &\longrightarrow \mathrm{image\ or\ belief},\\
\mathrm{belief} &\longrightarrow R_{\mathrm{psychological}}.
\end{align}

Here \(R_{\mathrm{psychological}}\) means psychological reality: the felt reality created or sustained by thought. The lecture is careful to distinguish this from nature. Thought did not create nature, though thought can describe nature, measure it, and make theories about it. But thought can create psychological realities, including illusions, and those illusions have real force in human life.

\begin{definition}
In this chapter, \(R_{\mathrm{psychological}}\) denotes a reality of belief, image, word, conclusion, or self-sense that has compelling force for the mind that holds it, whether or not it corresponds to what is true.
\end{definition}

This distinction matters because investigation stops when belief is treated as fact. If I say, ``This belief is reality,'' the door closes. The question is no longer alive. Therefore the investigation asks: can the mind see the belief as belief? Can it watch the movement by which the word becomes the thing?

\section{Can Thought Be Aware Of Itself?}

The final movement returns to correct action, but now with sharper tools. We have seen that authority is not the answer; that the me is disorder because it divides; that belief and word can become psychological reality; and that consciousness, at least as presently considered, is the proliferation of fragments.

So the question becomes:

\begin{equation}
Q_3:\qquad
\mathrm{Can\ consciousness\ be\ aware\ of\ itself?}
\end{equation}

Or, put still more simply:

\begin{equation}
T \overset{?}{\longrightarrow} \mathrm{awareness\ of\ }T,
\end{equation}

where \(T\) denotes thought as movement. Not thought as a fixed object, not a definition in a book, but thought moving: forming images, creating belief, dividing me from you, producing the structure of psychological reality.

\subsection{Question \& Answer}

\paragraph{Question.}
Can thought be aware of its own movement?

\paragraph{Answer.}
The lecture does not turn this into a method. It asks the listener to try it now. Can thought be aware of itself, its structure, its activity, what it has created, what it has done in the world? Can the doer be aware of the doing?

A tentative form is:

\begin{equation}
\mathrm{awareness}(T_{\mathrm{movement}})
\overset{?}{\longrightarrow}
\mathrm{radical\ change}.
\end{equation}

The dialogue briefly says, ``It stops,'' and then corrects the formulation: thought is undergoing a radical change. The correction is important. We should not make a doctrine out of the word ``stops.'' The point is not to assert a mechanism. The point is to stay with the question of whether observation changes the movement observed.

A physical analogy is mentioned: in observation, the object may not be fixed apart from the act of observation. The lecture does not develop this as physics, and we should not turn it into quantum mechanics. Its use here is analogical. The inward question is whether the movement of thought is altered when it is directly aware of itself.

\begin{figure}
\centering
\begin{tikzpicture}[
  node distance=0.85cm,
  box/.style={draw, rectangle, align=center, minimum width=0.78\linewidth, minimum height=0.82cm}
]
\node[box] (me) {the ``me'' as separation and disorder};
\node[box, below=of me] (c) {consciousness as collected fragments};
\node[box, below=of c] (t) {thought moving: word, image, belief, conclusion};
\node[box, below=of t] (aware) {can thought be aware of its own movement?};
\node[box, below=of aware] (change) {radical change?};
\draw[->] (me) -- (c);
\draw[->] (c) -- (t);
\draw[->] (t) -- (aware);
\draw[->] (aware) -- (change);
\end{tikzpicture}
\caption{The final inquiry is left open: awareness of thought's movement is not presented as a technique, but as the question to be tested.}
\end{figure}

The chapter should end where the lecture ends: not with consolation, but with a precise unfinished question. What is time in relation to this awareness? That question is deliberately held over.

\section{Summary}

The lecture moves from the public fact of human disorder to the intimate question of whether the root can change. It rejects secondary explanations without denying their partial truth. Habit, war, boredom, dependence, and cynicism all describe parts of the field, but none reaches the root.

The central discovery is that the one who wants to change may be part of the disorder. The planner creates new patterns while remaining old. This same structure appears as authority: disorder creates the demand for someone else to tell me what to do, and that demand sustains dependence.

From there the inquiry turns inward without becoming private self-improvement. Correct action cannot be supplied by authority, because action from disorder remains disorder. The mind must understand its own neuroticism: its fragmentation, its words, its beliefs, its psychological realities.

The final question is the most demanding one:

\begin{equation}
\mathrm{Can\ thought\ be\ aware\ of\ its\ own\ movement?}
\end{equation}

The lecture does not answer by prescribing a method. It leaves the question active, exact, and unfinished.